% English translation of checked_pages/esq_119_027.tex; source: cote 119, site 27, PDF page 28, typescript page 25.
% AI-translated by GPT-6 Astra (OpenAI, ChatGPT Pro mode).
\scanpage{119/27}
\mspage{25}

time seems ripe to rewrite a new version, in a modern
style, of Klein's classical book on the icosahedron and the other
Pythagorean polyhedra. Writing such an account of regular
2-polyhedra would be a splendid opportunity for a young researcher
to become familiar both with the geometry of polyhedra and
their connections with spherical, euclidean and hyperbolic geometries,
and with algebraic curves, and with the language and
basic techniques of modern algebraic geometry. Will one
ever be found to seize this opportunity ?

5. \qquad I would now like to say a few words about certain
reflections that led me to understand the need for new foundations
for ``geometric'' topology, in a direction quite
different from the notion of a topos, and independent even of the needs
of so-called ``abstract'' algebraic geometry (over general base fields
and rings). The initial problem, which began to
intrigue me what must already be some fifteen years ago, was
that of defining a theory of ``dévissage''\tnote{TN-027-01}{``Dévissage'' is retained, following the terminology standard. The source explains it here through reconstruction from canonical building pieces.}
of stratified structures, in order to reconstruct them by a canonical process
from ``building blocks'' canonically deduced from the given
structure. Probably the principal example that had then led me
to this question was the canonical stratification of a
singular algebraic variety (or a singular complex or real analytic
space) by the decreasing sequence of its successive
``singular loci''. But I probably already sensed the ubiquity
of stratified structures in practically every domain of
geometry (which others surely saw clearly long before
I did). Since then, I have seen such structures appear, in particular,
in every ``moduli'' situation for geometric objects susceptible
not only to continuous variation but, at the same time,
to phenomena of ``degeneracy'' (or ``specialization'') - the
strata then corresponding to the various ``levels of singularity''
(or to the associated combinatorial types) of the objects considered.

\margin{27}

